\section{Comparison with sign disorder}\label{sec:comparison}

We now transfer the Gaussian resolvent limits to sign couplings.  The
comparison uses the physical spin coordinates, so it requires no information
about eigenvectors for the sign model.  Its two inputs from
Section~\ref{sec:trace} are the joint Gaussian trace limit
\eqref{eq:source-40} and the vanishing quadratic-feature trace
\eqref{eq:source-43}.  The latter will control fourth derivatives of the former.
We use the Taylor replacement principle of \cite{Chatterjee2006}; its
application here rests on averaged derivative estimates for these dynamical
traces.

Let
\[
 \nu_t=(1-t)\mathsf N(0,1)+t\,\tfrac12(\delta_{-1}+\delta_1),
 \qquad 0\le t\le1,
\]
and write $\mathbb P_t,\mathbb E_t$ for independent disorder entries with
law $\nu_t$.  Thermal expectations continue to be denoted by $\mu$.
For a vector-valued feature $S$ define its resolvent trace by
\[
 B_S(z)=\langle S,(z+TH)^{-1}S\rangle_\mu,
 \qquad T=n^{2/3},\quad z>0,
\]
where the inner product sums the feature coordinates.  The two features used
below are
\[
 S^{[1]}(X)=n^{-1/3}X,\qquad
 S^{[2]}(X)=\frac{XX^{\mathsf t}}n.
\]
The second takes values in matrices with the Hilbert--Schmidt norm.  Set
$B_n^{[k]}(z)=B_{S^{[k]}}(z)$; thus $B_n^{[1]}(z)=B_n(z)$ by orthogonal
invariance of the sum defining the trace.  Positivity of $H$ gives the useful
deterministic bounds
\begin{equation}\label{eq:comparison-trivial-traces}
 0\le B_n^{[1]}(z)\le \frac{n^{1/3}}z,
 \qquad 0\le B_n^{[2]}(z)\le\frac1z.
\end{equation}

\begin{theorem}[Trace universality]\label{thm:comparison-universality}
For every $z>0$,
\[
 \sup_{0\le t\le1}\mathbb E_t B_n^{[2]}(z)\longrightarrow0.
\]
For each finite collection $z_1,\ldots,z_m>0$, the vector
$(B_n(z_1),\ldots,B_n(z_m))$ under sign disorder converges in law to the
Gaussian limiting vector in \eqref{eq:source-40}.  In particular the
coefficient $c_*$ in that limit is unchanged.
\end{theorem}

The proof first compares the bounded second-feature traces.  Their vanishing
then supplies the additional estimate needed to compare the first-feature
traces.  We begin with a static estimate that applies throughout the mixture.

\subsection{Uniform static overlap control}

For two independent spins under $\mu$, write
$\rho(X,X')=n^{-1}X\cdot X'$.  The following estimate also controls
correlations after multiplying $\mu$ by a density that depends on the disorder.

\begin{lemma}\label{lem:comparison-overlap}
For every $0<\epsilon\le1$ there is $a_\epsilon>0$ such that
\begin{equation}\label{eq:comparison-overlap-tail}
 \sup_{0\le t\le1}\mathbb P_t\left\{
   \mu^{\otimes2}(|\rho|\ge\epsilon)>e^{-a_\epsilon n}
 \right\}\longrightarrow0.
\end{equation}
Consequently, for fixed $C,K<\infty$, uniformly over nonnegative,
possibly disorder-dependent functions $p$ satisfying
$\mu p\le1$ and $\|p\|_\infty\le Cn^K$, the quantities
\[
 \frac1{\binom n2}\sum_{i<j}\big|\mu[pX_iX_j]\big|
\]
converge to zero in probability, uniformly in $t$.  They also converge to
zero in expectation, uniformly in $t$.
\end{lemma}

\begin{proof}
Let $Z$ denote the partition function with uniform cube probability as its
reference measure.  The moment generating function
$m_t(v)=\mathbb E_t e^{vJ_{12}}$ satisfies
\[
 m_t(v)\le e^{v^2/2},\qquad
 \log m_t(n^{-1/2})\ge\frac1{2n}-\frac C{n^2},
\]
with constants independent of $t$.  If $Z_A^{(2)}$ is the two-replica
partition function restricted to an event $A$ depending on $\rho$, independence
of the entries therefore gives
\begin{equation}\label{eq:comparison-annealed-pair}
 \frac{\mathbb E_t Z_A^{(2)}}{(\mathbb E_tZ)^2}
 \le C\,\mathbb E_{\mathrm{cube}}\!\left[
             e^{n\rho^2/2}\mathbf1_A\right].
\end{equation}
Here the expectation on the right is over two independent uniform cube spins.
Indeed the pair exponent uses
$(X_iX_j+X_i'X_j')^2=2+2X_iX_i'X_jX_j'$; summing over $i<j$ produces
$n\rho^2/2$ after division by the squared first moment, with only a bounded
factor remaining.

Put
\[
 I(r)=\frac{(1+r)\log(1+r)+(1-r)\log(1-r)}2,
 \qquad -1\le r\le1,
\]
with $0\log0=0$.  The binomial bound
$\mathbb P_{\mathrm{cube}}(\rho=r)\le e^{-nI(r)}$ and the fact that there
are $n+1$ possible overlaps show that the right side of
\eqref{eq:comparison-annealed-pair} is at most $C(n+1)$.  On
$A=\{|\rho|\ge\epsilon\}$ it is at most
\begin{equation}\label{eq:comparison-entropy-gap}
 C(n+1)e^{-n\kappa_\epsilon},\qquad
 \kappa_\epsilon=I(\epsilon)-\epsilon^2/2>0.
\end{equation}
For the strict inequality, $I'(r)-r=\operatorname{arctanh}(r)-r>0$ for
$0<r<1$, and continuity handles $r=1$.

We next obtain a lower bound for the quenched partition function.  The
unrestricted second-moment estimate and Paley--Zygmund imply
\[
 \mathbb P_t\{Z\ge\tfrac12\mathbb E_t Z\}\ge\frac c{n+1}.
\]
Clip every entry to $[-\log n,\log n]$, and let $Z_L$ be the resulting
partition function.  The probability that any entry changes is at most
$Cn^2e^{-(\log n)^2/2}$, uniformly in $t$.  Changing one clipped entry
changes $\log Z_L$ by at most $2\log n/\sqrt n$.  Bounded differences hence
gives, for each $\eta>0$,
\[
 \mathbb P_t\left\{
   |\log Z_L-\mathbb E_t\log Z_L|>\eta n
 \right\}
 \le2\exp\!\left(-\frac{c\eta^2n}{(\log n)^2}\right).
\]
This upper bound is smaller than $c/(2(n+1))$ for large $n$.  The preceding
Paley--Zygmund event thus forces
$\mathbb E_t\log Z_L\ge\log\mathbb E_tZ-\eta n$ for all sufficiently
large $n$, after adjusting $\eta$ by a fixed factor.  Concentration once more
shows that, for every fixed $\eta>0$,
\begin{equation}\label{eq:comparison-partition-lower}
 \sup_t\mathbb P_t\{
       Z<e^{-\eta n}\mathbb E_tZ\}\longrightarrow0.
\end{equation}
Use \eqref{eq:comparison-entropy-gap} and Markov's inequality to bound
$Z_A^{(2)}\le e^{-3\kappa_\epsilon n/4}(\mathbb E_tZ)^2$ with probability
tending to one.  On the event in \eqref{eq:comparison-partition-lower} with
$\eta=\kappa_\epsilon/8$, division by $Z^2$ proves
\eqref{eq:comparison-overlap-tail} with $a_\epsilon=\kappa_\epsilon/2$.

For the last assertion, an exact identity is
\begin{equation}\label{eq:comparison-weighted-overlap}
 \frac1{\binom n2}\sum_{i<j}\mu[pX_iX_j]^2
 =\frac{n\mu^{\otimes2}[p(X)p(X')\rho^2]-(\mu p)^2}{n-1}.
\end{equation}
On the event supplied by \eqref{eq:comparison-overlap-tail}, the replica
integral on the right is at most
$\epsilon^2+C^2n^{2K}e^{-a_\epsilon n}$.  First let $n\to\infty$ and then
$\epsilon\downarrow0$.  Cauchy--Schwarz gives the asserted average of
absolute correlations.  This average is bounded by $\mu p\le1$, which also
gives the uniform convergence in expectation.
\end{proof}

\subsection{Perturbing one interaction}

Fix an arbitrary disorder matrix, one of the two features $S$, and $z>0$.
All estimates in this subsection are deterministic unless an expectation
over disorder is displayed.  Let $e=\{i,j\}$ be an unordered pair, set
$q=X_iX_j$, and add $\theta$ to its interaction coefficient $J_{ij}/\sqrt n$.
Writing $u=\tanh\theta$, the perturbed Gibbs probability has density
$(1+uq)/(1+u\mu q)$ relative to $\mu$.

It is convenient to work first with the unnormalized measure
$(1+uq)\mu$.  Let $M(u)$ be the mass-plus-energy form for its heat-bath
operator:
\[
 M(u)(f,g)=z\mu[(1+uq)f\cdot g]
          +T\mathcal E_u^{\mathrm{un}}(f,g),
 \qquad b(u)(g)=\mu[(1+uq)S\cdot g].
\]
Here $\mathcal E_u^{\mathrm{un}}$ is the heat-bath Dirichlet form integrated
against $(1+uq)\mu$.  Put $M=M(0)$, write
$\|g\|_M^2=M(g,g)$ and $\|\cdot\|_*$ for its dual norm, and define
\[
 Y=M^{-1}b(0)=(z+TH)^{-1}S,\quad
 B=B_S(z)=M(Y,Y),\quad d=(1+B)^{1/2}.
\]
The feature coordinates are contracted in every form and dual pairing.
Expand $M(u)=\sum_{a\ge0}u^aM_a$ and $b(u)=b_0+ub_1$; the indices $a$
here denote Taylor coefficients, not sites.  The quantities that measure the
effect of the pair are
\begin{equation}\label{eq:comparison-sensitivities}
 \Gamma_e=\frac{T}{d^2}\sum_{k\in e}\mu[w_k|D_kY|^2],\qquad
 R_e=\frac1d\|b_1-M_1Y\|_*.
\end{equation}
We use $\operatorname{avg}_e$ for the average over the $\binom n2$ pairs.
Counting each site in $n-1$ pairs gives
\begin{equation}\label{eq:comparison-energy-average}
 0\le\Gamma_e\le1,\qquad
 \operatorname{avg}_e\Gamma_e
 =\frac2n\frac{T\mathcal E(Y)}{1+B}\le\frac2n.
\end{equation}

Gaussian and sign entries have the same first three moments, so the
comparison uses a fourth-order Taylor remainder.  The derivative estimates
below will require fourth-power averages of $R_e$.

\begin{lemma}[Residual estimates]\label{lem:comparison-residual}
Uniformly over the disorder and the pair $e$,
\begin{equation}\label{eq:source-44}
 R_e^2\le\frac{2T}{z}\Gamma_e.
\end{equation}
For $S=S^{[2]}$ one also has $R_e^2\le C_z$, whereas for $S=S^{[1]}$,
\begin{equation}\label{eq:source-45}
 R_e^2\le C_z\bigl(1+n^{1/3}\sqrt{B_n^{[2]}(z)}\bigr).
\end{equation}
Consequently $\operatorname{avg}_eR_e^2\le C_z n^{-1/3}$.  Its
fourth-power analogue is $\operatorname{avg}_eR_e^4\le C_z n^{-1/3}$ for
$S^{[2]}$, and
\begin{equation}\label{eq:source-46}
 \operatorname{avg}_eR_e^4
 \le C_z\bigl(n^{-1/3}+\sqrt{B_n^{[2]}(z)}\bigr)
 \qquad\text{for }S^{[1]}.
\end{equation}
\end{lemma}

\begin{proof}
For a site $k\notin e$, the energy multiplier in $M(u)$ is $1+uq$.
At an endpoint $k\in e$, let $s$ be the other spin in the pair.  Its
conditional mean changes from $t_k$ to
$(t_k+us)/(1+ust_k)$, while the marginal measure on $X_{-k}$ acquires
the factor $1+ust_k$.  Thus its energy term has the multiplier
\begin{equation}\label{eq:comparison-endpoint-multiplier}
 \frac{1-u^2}{1+ust_k}
\end{equation}
relative to $\mu[w_kD_kf\cdot D_kg]$.  Since $|t_k|\le1$, these formulas
show that $\|M_a\|_{M\to M^*}\le C^a$ and, for $a\ge2$, that $M_a$
uses only the two endpoint energies.  In particular,
\begin{equation}\label{eq:comparison-coefficient-locality}
 \begin{aligned}
 \|M_aY\|_*&\le C^a d\sqrt{\Gamma_e},&
 |M_a(Y,Y)|&\le C^a d^2\Gamma_e\quad(a\ge2),\\
 \|M_1Y\|_*&\le C\sqrt B.
 \end{aligned}
\end{equation}

Write $E_kg=\mu[g\mid X_{-k}]$.  Testing the equation $MY=b_0$ against
$qg$ gives the exact identity
\begin{equation}\label{eq:comparison-residual-identity}
 (b_1-M_1Y)(g)
  =T\sum_{k\in e}\mu[w_kD_kY\cdot sE_kg].
\end{equation}
To verify it, the mass terms and all sites outside $e$ cancel.  At an endpoint,
$D_k(qg)=s(g(X_k=1)+g(X_k=-1))/2$, whereas the first coefficient in
\eqref{eq:comparison-endpoint-multiplier} is $-st_k$.  The remaining sum is
$sE_kg$.  Cauchy--Schwarz and conditional Jensen now give
\[
 |(b_1-M_1Y)(g)|^2
 \le d^2\Gamma_e T\sum_{k\in e}\mu[w_k|E_kg|^2]
 \le\frac{2T}{z}d^2\Gamma_e\|g\|_M^2,
\]
which proves \eqref{eq:source-44}.

By \eqref{eq:comparison-coefficient-locality},
$R_e^2\le C_z(1+\|b_1\|_*^2)$.  For the second feature,
$\|qS^{[2]}\|_{L^2(\mu)}=1$, so $\|b_1\|_*^2\le1/z$.
For the first feature, use the stationary heat-bath process at the fixed
disorder to write
\begin{align*}
 \|b_1\|_*^2
 &=n^{1/3}\int_0^\infty e^{-zs}
   \mathbb E_\mu\!\left[
      q(X_0)q(X_{Ts})\rho(X_0,X_{Ts})\right]\,ds\\
 &\le n^{1/3}\left(\frac1z
      \int_0^\infty e^{-zs}
          \mathbb E_\mu[\rho(X_0,X_{Ts})^2]\,ds\right)^{1/2}
  =n^{1/3}\sqrt{B_n^{[2]}(z)/z}.
\end{align*}
The last equality is the Hilbert--Schmidt feature identity
$\langle XX^{\mathsf t}/n,X'X'^{\mathsf t}/n\rangle=\rho(X,X')^2$.
This proves \eqref{eq:source-45}.  Finally,
\eqref{eq:source-44} and \eqref{eq:comparison-energy-average} give the
second-power average.  Multiplying it by the respective pointwise bound on
$R_e^2$ proves both fourth-power estimates.
\end{proof}

The residual is small on average because a single pair touches only two
site energies.  The other contribution to a trace derivative is a correlation
with $q$.  The static overlap estimate controls precisely this contribution.
Define
\begin{equation}\label{eq:comparison-h}
 h_e=|\mu q|+d^{-2}|M_1(Y,Y)|+R_e+\sqrt{\Gamma_e}.
\end{equation}

\begin{lemma}[Averaged sensitivities]\label{lem:comparison-average}
For either feature, $\operatorname{avg}_e(h_e+h_e^4)\le C_z$
deterministically.  There is a deterministic sequence $\epsilon_n\to0$
such that, uniformly $t\in[0,1]$,
\begin{align}
 \mathbb E_t\operatorname{avg}_e(h_e+h_e^4)&\le\epsilon_n
       &&\text{for }S^{[2]},\label{eq:comparison-average-second}\\
 \mathbb E_t\operatorname{avg}_e(h_e+h_e^4)&\le\epsilon_n+
             C_z\mathbb E_t\sqrt{B_n^{[2]}(z)}
       &&\text{for }S^{[1]}.
       \label{eq:comparison-average-first}
\end{align}
\end{lemma}

\begin{proof}
The first two terms in \eqref{eq:comparison-h} are bounded by constants.
For the second of them introduce the nonnegative function
\[
 p_Y=\frac{z|Y|^2+T\sum_k w_k|D_kY|^2}{d^2}.
\]
Its mass is $B/(1+B)\le1$.  The Markov resolvent representation yields
$\|Y\|_\infty\le\|S\|_\infty/z$, with vector norm intended, and hence
$\|p_Y\|_\infty\le C_z n^2$ for the first feature and
$\|p_Y\|_\infty\le C_z n^{5/3}$ for the second.  The endpoint formula gives
\[
 \left|d^{-2}M_1(Y,Y)-\mu[qp_Y]\right|\le2\Gamma_e.
\]
Lemma~\ref{lem:comparison-overlap}, first with $p=1$ and then with $p=p_Y$,
therefore shows that these first two terms have vanishing pair averages in
expectation, uniformly in $t$.  Their fourth powers have the same property
because the terms are bounded.  The remaining terms are controlled by
\eqref{eq:comparison-energy-average},
Lemma~\ref{lem:comparison-residual}, and Cauchy--Schwarz:
\[
 \operatorname{avg}_e R_e\le C_z n^{-1/6},\qquad
 \operatorname{avg}_e\sqrt{\Gamma_e}\le\sqrt{2/n}.
\]
Combining these estimates proves \eqref{eq:comparison-average-second} and
\eqref{eq:comparison-average-first}.  The deterministic bound follows from
the same estimates and $B_n^{[2]}(z)\le1/z$.
\end{proof}

\subsection{Fourth derivatives of bounded trace observables}

The first trace need not be bounded uniformly in $n$.  We compare it through
the bounded coordinate $(1+B)^{-1}$.  The next lemma turns the preceding
sensitivities into the derivative estimate needed for replacement.

\begin{lemma}\label{lem:comparison-derivative}
Fix finitely many traces $B_a=B_n^{[k_a]}(z_a)$, with $k_a\in\{1,2\}$ and
$z_a>0$, and a smooth function $F$ on a neighborhood of $[0,1]^m$.
For an arbitrary base disorder and pair $e$, let $B_a(\theta)$ denote the
traces after adding $\theta$ to the interaction coefficient of $e$, and let
$\mathfrak h_e=\sum_{a=1}^m h_{e,a}$, with each sensitivity computed at the
base disorder.  Then for all sufficiently large $n$,
\begin{equation}\label{eq:comparison-fourth-derivative}
 \sup_{|\theta|\le2\log n/\sqrt n}
 \left|\frac{d^4}{d\theta^4}
 F\left(\frac1{1+B_1(\theta)},\ldots,
        \frac1{1+B_m(\theta)}\right)\right|
 \le C_F(\mathfrak h_e+\mathfrak h_e^4).
\end{equation}
At every real base disorder, derivatives of orders at most four also have
a deterministic bound polynomial in $n$.
\end{lemma}

\begin{proof}
We first consider one trace.  Complexify the finite-dimensional form space,
extending the forms and dual pairings complex bilinearly.
The coefficient bounds for $M(u)$ imply that
$M^{-1/2}M(u)M^{-1/2}$ differs from the identity in norm by $O(|u|)$ on a
fixed disk.  Its inverse is therefore analytic and uniformly bounded on a
smaller fixed disk.  Put $r(u)=b(u)-M(u)Y$ and
$A(u)=b(u)M(u)^{-1}b(u)$.  Completing the quadratic expression around $Y$
gives the algebraic identity
\[
 A(u)=2b(u)Y-M(u)(Y,Y)+r(u)M(u)^{-1}r(u).
\]
The coefficient of $u$ in $(A(u)-B)/d^2$ is
\[
 \frac{M_1(Y,Y)+2(b_1-M_1Y)(Y)}{d^2},
\]
whose absolute value is at most $Ch_e$.  For higher coefficients,
\eqref{eq:comparison-coefficient-locality} gives
\[
 \frac{\|r_1\|_*}{d}=R_e,\qquad
 \frac{\|r_a\|_*}{d}\le C^a\sqrt{\Gamma_e}\quad(a\ge2).
\]
The direct term $-M_a(Y,Y)/d^2$ is at most $C^a\Gamma_e$.  Coefficients of
the residual quadratic term are bounded, using the analytic inverse, by
$C^a(R_e^2+R_e\sqrt{\Gamma_e}+\Gamma_e)\le C^ah_e^2$.
Thus every coefficient of order $a\ge2$ in $(A(u)-B)/d^2$ is bounded by
$C^ah_e^2$.

The actual trace divides $A(u)$ by the total mass $1+u\mu q$.  Thus, with
\[
 \delta(u)=\frac{B(u)-B}{d^2},
\]
its linear coefficient is $O(h_e)$ and its coefficients of order $a\ge2$
are bounded by $C^a(h_e+h_e^2)$.  From
\eqref{eq:source-45} and \eqref{eq:comparison-trivial-traces},
$h_e\le C_z(1+n^{1/6})$ for the first feature; for the second it is bounded.
Consequently $|u|h_e=o(1)$ uniformly for $|u|\le2\log n/\sqrt n$.
Summing the coefficient bounds on this interval gives
\[
 \sup|\delta(u)|=o(1),\qquad
 |\delta'(u)|\le Ch_e,\qquad
 |\delta^{(a)}(u)|\le C(h_e+h_e^2),\quad 2\le a\le4.
\]
Since
\[
 \frac1{1+B(u)}=\frac1{d^2}\frac1{1+\delta(u)},
\]
the chain rule bounds its fourth derivative, and that of any fixed smooth
test of it, by $C(h_e+h_e^4)$.  The derivatives of $u=\tanh\theta$ through
order four are bounded on the interval in question, so the same estimate
holds for interaction displacement $\theta$.

For several traces apply the multivariable chain rule and bound every mixed
product using $\mathfrak h_e+\mathfrak h_e^4$.  Finally, at any real base
point the coefficient argument applies at $\theta=0$, where the denominator
of the compactification is at least one.  The bound
$h_e\le C(1+n^{1/6})$ supplies the claimed polynomial bounds, uniformly in
that base point.
\end{proof}

\subsection{Replacement and the two-feature argument}

We have obtained a small average fourth derivative under a complete iid
mixture disorder.  Taylor replacement removes an entry before expanding,
so its expansion point does not itself have that law.  The following
independent completion lets us use the averaged estimates without changing
the Taylor coefficients.

Let
\[
 \Phi(J)=F\left((1+B_n^{[k_1]}(z_1))^{-1},\ldots,
               (1+B_n^{[k_m]}(z_m))^{-1}\right).
\]
Differentiating the finite product mixture measure gives exactly
\begin{equation}\label{eq:comparison-mixture-derivative}
 \frac d{dt}\mathbb E_t\Phi
  =\sum_e\mathbb E_{t,-e}\!\left[
       \mathbb E_{\xi}\Phi(J_{-e},\xi)
       -\mathbb E_g\Phi(J_{-e},g)\right],
\end{equation}
where $\xi$ is a symmetric sign, $g$ is standard normal, and
$\mathbb E_{t,-e}$ averages all entries other than $e$.
For fixed $J_{-e}$, Taylor-expand $x\mapsto\Phi(J_{-e},x)$ at zero through
degree three.  These polynomial terms cancel in
\eqref{eq:comparison-mixture-derivative}, since $\xi$ and $g$ have the same
first three moments.

Independently of both proposed values, draw $J_e^0$ with law $\nu_t$, and
write $J^0=(J_{-e},J_e^0)$.  If $|x|,|J_e^0|\le\log n$, every point of the
Taylor segment between $0$ and $x$ differs from $J_e^0$ by at most
$2\log n$.  Lemma~\ref{lem:comparison-derivative}, with $J^0$ as its base,
therefore bounds the Taylor remainder by
\begin{equation}\label{eq:comparison-remainder}
 C_F\frac{|x|^4}{n^2}
       \bigl(\mathfrak h_e(J^0)+\mathfrak h_e(J^0)^4\bigr).
\end{equation}
The factor $n^{-2}$ changes four derivatives in the interaction coefficient
$J_e/\sqrt n$ into four derivatives in $J_e$.  The factor $|x|^4$ is measured
from the Taylor center $0$.  The independently completed entry is used only
to bound the derivatives along the segment, so it does not enter this factor.

The events $|g|>\log n$ and $|J_e^0|>\log n$ have probability smaller than
every fixed inverse power of $n$, uniformly in $t$.  The boundedness of
$\Phi$, together with the polynomial derivative bounds in
Lemma~\ref{lem:comparison-derivative}, bounds the Taylor polynomial and its
remainder by a polynomial in $n$ times $1+|x|^4$.  Gaussian tail moments thus
make the omitted contributions $o(n^{-2})$, uniformly in $J_{-e}$ and $t$.
On the remaining event, enlarge the right side of
\eqref{eq:comparison-remainder} by removing its indicator.  Its $|x|^4$
factor is independent of the completed base $J^0$ and has bounded expectation.
Moreover, if $e$ is chosen uniformly, $J^0$ has the full iid mixture law
independently of $e$.  Summing the remainders in
\eqref{eq:comparison-mixture-derivative} therefore yields
\begin{equation}\label{eq:comparison-master-bound}
 \left|\frac d{dt}\mathbb E_t\Phi\right|
 \le C_F\mathbb E_t\operatorname{avg}_e
             (\mathfrak h_e+\mathfrak h_e^4)+o(1),
\end{equation}
uniformly for $0\le t\le1$.

\begin{proof}[Proof of Theorem~\ref{thm:comparison-universality}]
First use only second-feature traces in \eqref{eq:comparison-master-bound}.
For finitely many traces,
$\mathfrak h_e+\mathfrak h_e^4\le C_m\sum_a(h_{e,a}+h_{e,a}^4)$.
Equation~\eqref{eq:comparison-average-second} therefore makes the mixture
derivative tend to zero uniformly in $t$.  In particular, for one $z$, choose
the test $F(v)=1-v$.  The Gaussian convergence
\eqref{eq:source-43} and the bound $B_n^{[2]}(z)\le1/z$ imply
\[
 \sup_t\mathbb E_t\frac{B_n^{[2]}(z)}{1+B_n^{[2]}(z)}\longrightarrow0.
\]
Since $b\le(1+1/z)b/(1+b)$ for $0\le b\le1/z$, this proves the first
assertion of the theorem and, by Jensen,
$\sup_t\mathbb E_t\sqrt{B_n^{[2]}(z)}\to0$.

Now use first-feature traces.  Equation~\eqref{eq:comparison-average-first}
makes their right side in \eqref{eq:comparison-master-bound} tend to zero
uniformly in $t$.  Integrating from $t=0$ to $t=1$ shows equality of the
limiting expectations of every smooth test of finitely many compactified
traces.  Such tests determine probability laws on $[0,1]^m$.  By
Proposition~\ref{prop:trace-goe}, the Gaussian trace vector converges to the
finite vector in \eqref{eq:source-40}; its compactification therefore has
all coordinates strictly positive almost surely.  The continuous mapping
theorem for the inverse map $v\mapsto v^{-1}-1$ recovers the same joint
limit for the sign traces.
\end{proof}
